\documentclass[11pt]{article}
\usepackage[margin=1.1in]{geometry}
\usepackage{amsmath,amssymb}
\usepackage{booktabs}
\usepackage{graphicx}
\usepackage[colorlinks=true,linkcolor=blue,citecolor=blue,urlcolor=blue]{hyperref}

\newcommand{\rr}{\mathbf{r}}
\newcommand{\kk}{\mathbf{k}}
\newcommand{\EE}{\mathbf{E}}
\newcommand{\HH}{\mathbf{H}}
\newcommand{\Mv}[1]{\mathbf{M}^{(#1)}}
\newcommand{\Nv}[1]{\mathbf{N}^{(#1)}}
\newcommand{\rhat}{\hat{\mathbf{r}}}
\newcommand{\thhat}{\hat{\boldsymbol{\theta}}}
\newcommand{\phhat}{\hat{\boldsymbol{\varphi}}}
\newcommand{\khat}{\hat{\mathbf{k}}}
\newcommand{\ehat}{\hat{\mathbf{e}}}
\newcommand{\Ptil}{\widetilde{P}}

\title{Operational Manual: T-Matrix Extraction from CST Studio Suite\\
Frequency-Domain Simulations}
\author{\texttt{cst\_tmatrix} package documentation}
\date{\today}

\begin{document}
\maketitle

\begin{abstract}
This manual documents the theory, software architecture, operating
procedure, and validation of the \texttt{cst\_tmatrix} package, which
computes transition matrices (T-matrices) of arbitrary, spatially
localized scatterers from full-wave simulations performed with the
frequency-domain (FDFD) solver of CST Studio Suite.  The incident and
scattered multipole amplitudes are both obtained from the simulated total
field on a spherical surface, which makes the extraction independent of
the amplitude and phase conventions of the excitation.  Results are
stored in the community \texttt{tmat.h5} data format
\cite{Asadova2025format} used by the multiple-scattering code
\emph{treams} \cite{Beutel2024treams}.  The method follows the
formulation of the Neural Analytic Modal Method (NAMM) document
\cite{NAMM}, with the conventions closed and verified as described here.
\end{abstract}

\tableofcontents

%======================================================================
\section{Introduction}
The transition matrix (T-matrix) of a localized object relates the
coefficients of an incident electromagnetic field, expanded in regular
vector spherical waves, to the coefficients of the scattered field,
expanded in outgoing vector spherical waves
\cite{Waterman1965,Mishchenko2002}.  It encodes the complete linear
electromagnetic response of the object at a given frequency and is the
natural building block for multiple-scattering computations of clusters,
metasurfaces, and disordered media through the Foldy--Lax equations and
translation-addition theorems \cite{Foldy1945,Tsang2001}.

For objects without analytic solutions the T-matrix must be computed
numerically.  The approach implemented here illuminates the object with a
set of plane waves from distributed directions, solves each scattering
problem with the frequency-domain solver of CST Studio
Suite, records the complex electric and magnetic fields on a spherical
surface enclosing the object, and determines the T-matrix from the
least-squares relation between the incident and scattered expansion
coefficients \cite{Fruhnert2017}.  Two properties distinguish the present
implementation:
\begin{enumerate}
\item \emph{Self-calibrating field separation.}  The incident and
scattered coefficients are both extracted from the same recorded total
field (Sec.~\ref{sec:separation}); no assumption about the source
amplitude or phase enters the T-matrix.
\item \emph{Verified conventions.}  Every sign, phase, and normalization
choice is fixed either by an offline test against an independent
reference or by comparison with a full-wave solution of a problem with a
known answer (Sec.~\ref{sec:validation}).
\end{enumerate}

%======================================================================
\section{Theoretical formulation}

\subsection{Conventions and vector spherical waves}
\label{sec:conventions}
All internal computations use the engineering time convention
$e^{+j\omega t}$, consistent with CST.  Outgoing waves therefore carry
the spherical Hankel function of the second kind,
$h_l^{(2)}(kr)$; regular waves carry the spherical Bessel function
$j_l(kr)$.  Spherical harmonics are orthonormal with the
Condon--Shortley phase,
\begin{equation}
Y_{lm}(\theta,\varphi) = \Ptil_l^m(\cos\theta)\, e^{jm\varphi},
\qquad
\int_{4\pi} Y_{lm} Y_{l'm'}^* \, d\Omega = \delta_{ll'}\delta_{mm'},
\end{equation}
where $\Ptil_l^m$ are the fully normalized associated Legendre functions
(the normalization of Jackson \cite{Jackson1999}).  With
$\gamma_l = 1/\sqrt{l(l+1)}$ and the angular functions
\begin{equation}
\pi_{lm}(\theta) = \frac{m\,\Ptil_l^m(\cos\theta)}{\sin\theta},
\qquad
\tau_{lm}(\theta) = \frac{d \Ptil_l^m(\cos\theta)}{d\theta},
\label{eq:pitau}
\end{equation}
the vector spherical waves are, for radial kind $c$ ($c=1$: $z_l = j_l$;
$c=3$: $z_l = h_l^{(2)}$),
\begin{align}
\Mv{c}_{lm}(\rr) &= \gamma_l\, z_l(kr)
 \left[\, j\pi_{lm}\,\thhat - \tau_{lm}\,\phhat \,\right] e^{jm\varphi},
\label{eq:Mdef}\\
\Nv{c}_{lm}(\rr) &= \gamma_l \left[\, l(l+1)\,\frac{z_l(kr)}{kr}\,
 \Ptil_l^m\, e^{jm\varphi}\,\rhat
 + \frac{[kr\,z_l(kr)]'}{kr}
 \left(\tau_{lm}\,\thhat + j\pi_{lm}\,\phhat\right) e^{jm\varphi} \right],
\label{eq:Ndef}
\end{align}
with $\Nv{c} = k^{-1}\nabla\times\Mv{c}$.  The tangential vector
harmonics obtained from Eqs.~(\ref{eq:Mdef})--(\ref{eq:Ndef}) are
orthonormal over the sphere.  Fields are expanded as
\begin{equation}
\EE_{\mathrm{inc}} = \sum_{l=1}^{l_{\max}}\sum_{m=-l}^{l}
 a^M_{lm}\Mv{1}_{lm} + a^N_{lm}\Nv{1}_{lm},
\qquad
\EE_{\mathrm{sca}} = \sum_{l,m}
 f^M_{lm}\Mv{3}_{lm} + f^N_{lm}\Nv{3}_{lm},
\label{eq:expansions}
\end{equation}
and the magnetic field follows from the curl relations as
$Z_0\HH = j\sum a^M \Nv{} + a^N \Mv{}$ (same radial kind as the
corresponding electric term), with $Z_0$ the free-space impedance.
Coefficients are stacked as $[\,M\text{ block};\,N\text{ block}\,]$ with
$l = 1\ldots l_{\max}$ outer and $m = -l \ldots l$ inner, giving
$N = 2\,l_{\max}(l_{\max}+2)$ modes, and the T-matrix is defined by
\begin{equation}
\mathbf{f} = \mathsf{T}\,\mathbf{a}.
\label{eq:Tdef}
\end{equation}

\subsection{Plane-wave expansion}
A unit-amplitude linear plane wave propagating along
$\khat(\theta_i,\varphi_i)$ with polarization vector $\ehat$ equal to the
spherical unit vector $\thhat$ or $\phhat$ at the propagation direction,
\begin{equation}
\EE(\rr) = \ehat\, e^{-jk\,\khat\cdot\rr},
\end{equation}
has the regular-wave coefficients
\begin{equation}
a^M_{lm} = 4\pi(-j)^l\; \ehat\cdot\mathbf{C}^*_{lm}(\khat),
\qquad
a^N_{lm} = 4\pi(-j)^{l-1}\; \ehat\cdot\mathbf{B}^*_{lm}(\khat),
\label{eq:pwcoeffs}
\end{equation}
with the tangential harmonics
$\mathbf{C}_{lm} = \gamma_l (j\pi_{lm}\thhat - \tau_{lm}\phhat)
 e^{jm\varphi}$ and
$\mathbf{B}_{lm} = \gamma_l (\tau_{lm}\thhat + j\pi_{lm}\phhat)
 e^{jm\varphi}$ evaluated at the propagation direction.
Equation~(\ref{eq:pwcoeffs}) was derived by numerical projection onto the
basis of Sec.~\ref{sec:conventions} rather than transcribed from the
literature, and is fixed by a test that reconstructs the plane-wave field
from the coefficients to a relative accuracy of $10^{-13}$.

\subsection{Multipole truncation}
\label{sec:truncation}
Wiscombe's criterion \cite{Wiscombe1980}, applied to the size parameter
$x = k_{\max} a$ of the circumscribing sphere of radius $a$,
\begin{equation}
w(x) = \left\lceil x + 4x^{1/3} + 1 \right\rceil \;(x\le 8),
\qquad
w(x) = \left\lceil x + 4.05\,x^{1/3} + 2 \right\rceil \;(8<x<4200),
\label{eq:wiscombe}
\end{equation}
is an upper bound rather than a prescription.  It gives the order at
which the Mie series has converged to machine precision, which is far
smaller than the uncertainty of the present extraction
($10^{-3}$--$10^{-2}$ in the coefficients;
Sec.~\ref{sec:failuremodes}).  Multipole orders whose true amplitudes
lie below that uncertainty cannot be determined from the data; retaining
them adds noise to the matrix and increases the cost twice over, since
$l_{\max}$ fixes both the number of modes
$N = 2\,l_{\max}(l_{\max}+2)$ --- and with it the number of solver runs
--- and the minimum monitor radius (Sec.~\ref{sec:monitorplacement}),
and therefore the size of the simulation domain.

In practice one takes Eq.~(\ref{eq:wiscombe}) as the upper bound,
determines the per-order amplitudes $|f_l|$ of the scattered field from
a single test solve (or from a previous extraction of a similar object),
finds the largest order still resolved above the uncertainty, and sets
$l_{\max}$ one order higher.  Physical and spurious content are easy to
tell apart in a stored matrix: physical amplitudes decrease with $l$ and
depend on frequency, whereas noise amplified by the ill-conditioning at
high orders \emph{increases} with $l$ and is essentially independent of
frequency.  As an illustration, a gold spoke-and-wheel resonator
extracted at the Wiscombe value $l_{\max}=9$ had $81\%$ of its matrix
norm in the $l=9$ block, constant to four digits from 10 to 34~THz
--- clearly not a property of the scatterer; the
measurable content ended near $l = 4$, and re-extraction at
$l_{\max}=5$ satisfied all diagnostics.  The per-frequency Wiscombe
order and the fraction of the matrix norm above it are stored with
every extraction (\texttt{noise\_order\_content},
Sec.~\ref{sec:diagnostics}); a masking utility sets the undetermined
orders to zero before they enter sums over modes.

\subsection{Self-calibrating field separation}
\label{sec:separation}
The field monitor records the \emph{total} field
$\EE = \EE_{\mathrm{inc}} + \EE_{\mathrm{sca}}$ on a sphere of radius
$R$, chosen for each frequency sub-range by the placement criteria of
Sec.~\ref{sec:monitorplacement}.  Projecting the tangential components of
$\EE$ and $Z_0\HH$ onto $\mathbf{C}_{lm}$ and $\mathbf{B}_{lm}$ yields
four scalars per mode,
\begin{equation}
p^C_{lm} = \langle \mathbf{C}_{lm}, \EE_t\rangle,\quad
p^B_{lm} = \langle \mathbf{B}_{lm}, \EE_t\rangle,\quad
h^C_{lm} = \langle \mathbf{C}_{lm}, Z_0\HH_t\rangle,\quad
h^B_{lm} = \langle \mathbf{B}_{lm}, Z_0\HH_t\rangle,
\end{equation}
which are related to the unknown incident and scattered coefficients by
two $2\times 2$ linear systems per mode,
\begin{equation}
\begin{pmatrix} p^C_{lm}\\ h^B_{lm} \end{pmatrix}
=
\begin{pmatrix} j_l & h^{(2)}_l\\[2pt]
 j\,\psi_l' /kr\;\; & j\,\xi_l'/kr \end{pmatrix}
\begin{pmatrix} a^M_{lm}\\ f^M_{lm} \end{pmatrix},
\qquad
\begin{pmatrix} p^B_{lm}\\ h^C_{lm} \end{pmatrix}
=
\begin{pmatrix} \psi_l'/kr & \xi_l'/kr\\[2pt]
 j\,j_l & j\,h^{(2)}_l \end{pmatrix}
\begin{pmatrix} a^N_{lm}\\ f^N_{lm} \end{pmatrix},
\label{eq:separation}
\end{equation}
where $\psi_l = kr\,j_l(kr)$ and $\xi_l = kr\,h^{(2)}_l(kr)$ are the
Riccati--Bessel functions evaluated at $kr = kR$.  The determinants of
the two systems in Eq.~(\ref{eq:separation}) are proportional to the
Wronskian $W\{j_l, h^{(2)}_l\} \neq 0$, so the decomposition of the
total field into incident and scattered parts always has a unique
solution.  Its sensitivity to errors in the field data is another
matter: for $kR < l$ the Bessel function $j_l(kR)$ is exponentially
small (the classically forbidden region below the centrifugal barrier),
the matrix in Eq.~(\ref{eq:separation}) becomes nearly singular, and
small errors in the recorded field produce large errors in $a_{lm}$ at
that order.  This is the origin of the first placement criterion in
Sec.~\ref{sec:monitorplacement}.  Because $\mathbf{a}$ and $\mathbf{f}$
are obtained from the same data, any amplitude or phase convention of
the source cancels exactly in Eq.~(\ref{eq:lstsq}) below.  CST documents
its plane-wave source as unit-amplitude with phase reference at the
global origin; the extraction confirms this
(Sec.~\ref{sec:validation}) but does not rely on it.

Since the incident field is known analytically, comparing the measured
coefficients $\mathbf{a}$ with Eq.~(\ref{eq:pwcoeffs}) provides a
diagnostic of the fidelity of the recorded field itself --- the
\emph{incident deviation}, stored per frequency with every extraction.
It is nonzero because of the discretization error of the solver and
grows when the recorded field is degraded for any reason.  The two
mechanisms observed in practice are a monitor sphere placed inside the
reactive near field of the scatterer's higher multipoles
(Sec.~\ref{sec:monitorplacement}), which raises the deviation to
$10^{-1}$, and insufficient mesh resolution at the upper end of the
frequency range,
which raises it gradually to a few times $10^{-2}$.  In template mode
(Sec.~\ref{sec:template}), where the geometry is supplied by the user
rather than generated by the package, a large deviation can also
indicate a geometry that is not centered on the origin or a project
built in different units than declared; the programmatic path excludes
these by construction.

\subsection{Monitor placement}
\label{sec:monitorplacement}
The monitor radius $R$ must satisfy two conditions.  Both are most
restrictive at the lowest frequency of the extraction range, where $kR$
is smallest:
\begin{align}
\text{(i)}\quad
 & k_{\min} R \;\ge\; l_{\max} + 2,
\label{eq:critcond}\\
\text{(ii)}\quad
 & k R \;\ge\; w\!\left(k a\right) + 1 \quad \text{at every frequency of
 the range},
\label{eq:critnear}
\end{align}
with $w$ the Wiscombe order of Eq.~(\ref{eq:wiscombe}) and $a$ the
circumscribing radius.  Condition (i) keeps the incident/scattered
decomposition of Sec.~\ref{sec:separation} well conditioned at
$l = l_{\max}$: for $kR < l$ the matrix elements containing $j_l(kR)$
are exponentially small and the inversion amplifies field errors.
Condition (ii) is independent of the truncation.  The scatterer
radiates multipoles up to order $w(ka)$ whether or not the projection
retains them, and orders with $l > kR$ are still evanescent at the
monitor radius: their amplitudes grow steeply toward the scatterer,
they enlarge the total field recorded on the sphere, and the solver's
relative error on that enlarged field enters every projected
coefficient as an absolute error.  The left side of (ii) grows with
$k$ faster than $w(ka)$ does, so the condition need only be checked at
$k_{\min}$.

Both conditions, and the $+1$ in (ii), come from observed failures.
For (i): an extraction with $k_{\min}R = 0.90$ and $l_{\max}=3$ gave
incident deviations up to $0.99$, concentrated at $l = l_{\max}$.  For
(ii): five extractions, all satisfying (i), are summarized by the
continuous quantity $n(x) = x + 4x^{1/3} + 1$ (whose integer ceiling is
$w$).  Extractions with $k_{\min}R - n = +0.93$, $+0.75$, and $+0.67$
gave incident deviations below $10^{-2}$; those with $+0.44$ and
$-0.25$ gave deviations of $0.11$--$0.15$, singular values of
$\mathsf{S}$ up to $1.37$, and an $l=l_{\max}$ block roughly an order
of magnitude too large, largest at the low-frequency end of the range
--- and, where two sub-ranges adjoin, discontinuous at their
junction, which is what distinguishes this mechanism from insufficient
mesh resolution (Sec.~\ref{sec:failuremodes}).  The integer ceiling
initially obscured the distinction: two of these extractions both had
$kR = w = 8$ while differing by $0.31$ in $kR - n$, and a rule
requiring only $kR \ge w$, consistent with the first four data points,
failed on the fifth.  Requiring $kR \ge w + 1$ places every future
extraction at $kR - n \ge 1$, above all observations that satisfied the
diagnostics.  \texttt{recommended\_monitor\_factor} computes the
smallest $R$ satisfying both conditions and is the default
($\texttt{monitor\_factor=None}$); a manually supplied factor is
checked at extraction start.  Note also the coupling to
Sec.~\ref{sec:truncation}: condition (i) ties $R$ to $l_{\max}$, so an
unnecessarily large truncation enlarges the simulation domain and its
memory use roughly as the cube of $R$ --- a further reason to choose
$l_{\max}$ from the measured multipole content rather than from
Eq.~(\ref{eq:wiscombe}) alone.

\subsection{Least-squares extraction and illumination set}
\label{sec:lstsq}
Collecting the coefficient vectors of $K$ illuminations as columns of
$N\times K$ matrices $\mathsf{A}$ and $\mathsf{F}$, the T-matrix is the
least-squares solution of
\begin{equation}
\mathsf{F} = \mathsf{T}\,\mathsf{A}
\qquad\Longrightarrow\qquad
\mathsf{T} = \mathsf{F}\mathsf{A}^{\dagger}
 \left(\mathsf{A}\mathsf{A}^{\dagger}\right)^{-1},
\label{eq:lstsq}
\end{equation}
computed by a rank-revealing solver.  The illumination set consists of
$K/2$ directions with two orthogonal polarizations each.  Two direction
schemes are provided:
\begin{description}
\item[Fibonacci (default).]  Quasi-uniform spiral points, available for
any direction count and free of the coordinate poles.  The number of
directions is $\lceil \eta N/2 \rceil$ with overdetermination factor
$\eta$ (default $1.5$); the observed condition number of $\mathsf{A}$ is
close to unity.
\item[Lebedev.]  Nodes of Lebedev quadrature rules \cite{Lebedev1976}.
Because the least-squares formulation does not integrate over directions,
the quadrature exactness of the rule is not required here; the practical
difference from the Fibonacci scheme is that Lebedev rules exist only at
specific point counts, so the number of solver runs is rounded up to the
next available rule.  Rules with 6, 14, and 26 points are built in with
exact weights; larger standard rules can be supplied as data files and
are verified by integrating spherical harmonics before use.
\end{description}
Overdetermination averages solver noise: in synthetic tests, additive
field noise of $10^{-3}$ produces a T-matrix error of
$4\times 10^{-5}$.  The computation consists of
$K = 2\lceil \eta N/2\rceil$ solver runs, each of which
covers \emph{all} requested frequencies through field monitors placed
at exact solver samples (Table~\ref{tab:cost}); symmetry reduction
(Sec.~\ref{sec:symmetryreduction}) divides $K$ by up to the group
order.  The solver treats each frequency sample as a separate linear
system, so the total time scales as the \emph{product} of the
illumination count and the number of frequency samples, plus a meshing
overhead per run (CST regenerates the mesh on every solver start;
Appendix~\ref{app:cstbugs}).  Reducing either factor helps
proportionally; once the illumination count has been reduced by
symmetry, the frequency grid is the remaining large factor.

\begin{table}[ht]
\centering
\caption{Number of frequency-domain solver runs for overdetermination
factor $\eta = 1.5$.}
\label{tab:cost}
\begin{tabular}{cccc}
\toprule
$l_{\max}$ & modes $N$ & directions & solver runs\\
\midrule
2 & 16 & 12 & 24\\
3 & 30 & 23 & 46\\
5 & 70 & 53 & 106\\
8 & 160 & 120 & 240\\
\bottomrule
\end{tabular}
\end{table}

\subsection{Symmetry-reduced illumination}
\label{sec:symmetryreduction}
If the scatterer is invariant under a point group $G$ (axial groups
$C_n$ and $C_{nv}$ about $z$ are currently supported), the T-matrix
commutes with the representation $D(g)$ of every group element on the
mode basis \cite{Waterman1971}.  Each solved illumination
$(\mathbf{a}, \mathbf{f})$ therefore yields $|G|$ columns
$(D(g)\mathbf{a}, D(g)\mathbf{f})$, of which only one requires a solver
run; for $C_{4v}$ ($|G| = 8$) the measured saving is a factor of
$7.2$--$7.8$ in solver runs.  The reduction is enabled through
\texttt{ExtractionConfig(symmetry\_n\_fold,
symmetry\_mirror\_phi0\_deg)}.  The assumed symmetry should be verified
beforehand: geometrically, with the STL-based point-group detector
(\texttt{postprocess.symmetry.detect\_axial\_point\_group}), and, when
a previous extraction of the object exists, by computing
$\|\mathsf{T}D - D\mathsf{T}\|/\|\mathsf{T}\|$ for the candidate
operations --- if this is comparable to the reciprocity deviation of
the same matrix, the symmetry holds to within the accuracy of the
data, since reciprocity sets the scale of the extraction error.

Three details of the implementation deserve comment, because each
involves a sign or ordering convention that a test suite sharing the
same convention cannot detect; all three were fixed by comparison with
an independently computed quantity.
\begin{enumerate}
\item Rotating the illumination direction corresponds to the
\emph{inverse} rotation acting on the coefficients:
$\mathbf{a}(\theta, \varphi + \alpha) = D(-\alpha)\,
\mathbf{a}(\theta, \varphi)$.  The sign of $\alpha$ cannot be checked
by internal consistency --- tests constructed with either sign are
consistent with themselves --- so it was fixed by evaluating the
plane-wave expansion, Eq.~(\ref{eq:pwcoeffs}), directly at rotated
angles and comparing (agreement to $3\times 10^{-16}$ for the correct
sign; order unity for the wrong one).
\item A vertical mirror plane at azimuth $\varphi_0$ acts as
$D(\sigma_{\varphi_0})\mathbf{a}(\theta,\varphi,\hat{e}) = s\,
\mathbf{a}(\theta, 2\varphi_0 - \varphi, \hat{e})$ with $s = +1$ for
$\thhat$-polarization and $s = -1$ for $\phhat$-polarization.  The
matrix $D(\sigma)$ is obtained numerically, by evaluating each basis
mode at the mirrored points, applying the polar-/axial-vector
transformation to $(\EE, \HH)$, and projecting back onto the basis;
the construction is checked through the involution $D^2 = I$.
\item The solved directions must be spread over the sphere, e.g.\ by
farthest-point selection.  The group operations preserve the polar
angle $\theta$, so the orbit of a direction adds only azimuthal
copies.  Taking candidate directions in the order of the Fibonacci
spiral, which begins near the pole, concentrates all solved directions
in a polar cap; the resulting system of columns had numerical rank
$38$ of $48$ and condition number $4\times 10^5$, and the matrix
obtained from it violated reciprocity and passivity at order unity on
measured data even though the same construction reproduced synthetic
(noise-free) matrices to $10^{-8}$ --- a reminder that recovery tests
without noise say nothing about conditioning.  Farthest-point
selection of the same number of directions gives condition numbers of
$10$--$25$.
\end{enumerate}
Illuminations on the symmetry axis or in a mirror plane have smaller
orbits (some group elements map them to themselves) and are excluded
automatically.  The expansion is applied before the least-squares
solve, so matrix elements forbidden by the enforced group vanish
identically in the result.  Elements forbidden by the full symmetry of
the object but not by the enforced subgroup --- for a body of
revolution under $C_{4v}$, those with $m - m' = \pm 4$ --- are then a
direct measurement of the extraction error, since their true values
are zero.

This observation extends to a post-processing step
(\texttt{postprocess.symmetry.project\_onto\_symmetry}).  For a
symmetry that has been verified independently, the true T-matrix is
invariant under the group, so the non-invariant part of a measured
matrix is extraction error; its relative magnitude,
$\|\mathsf{T} - P(\mathsf{T})\|_F/\|\mathsf{T}\|_F$ with $P$ the group
average, quantifies that error using every matrix element, and
$P(\mathsf{T})$ is the corresponding cleaned matrix.  The group
average is preferable to setting only the forbidden elements to zero:
it is the orthogonal projection onto the invariant subspace, so it
also averages symmetry-related allowed elements and reduces their
error.  Two rules govern its use.  The projection may only be applied
for a verified symmetry --- applied to a symmetry the object lacks, it
replaces physics with an assumption.  And the raw matrix should be
kept on disk with the removed fraction recorded, the projection being
applied downstream: the diagnostics of the raw matrix must be examined
first, because an error mechanism with a symmetric component (a
misplaced monitor, for instance) is only partially removed by $P$ and
the projected matrix then appears more accurate than the underlying
data are.

\subsection{Surface quadrature}
The projections in Eq.~(\ref{eq:separation}) are evaluated with a product
quadrature: Gauss--Legendre nodes in $\cos\theta$ and uniform nodes in
$\varphi$.  This rule integrates spherical harmonics exactly up to a
designed degree, like a Lebedev rule, but is available at any order.  The
design degree accounts for the full angular bandwidth of the recorded
field: the incident plane wave carries angular content up to
approximately $kR$ on the monitor sphere, so the quadrature is sized for
degree $l_{\mathrm{field}} + l_{\max}$ with
$l_{\mathrm{field}} = \lceil kR + 4.05(kR)^{1/3} + 2\rceil$ plus a safety
margin.  Undersizing this quadrature aliases incident-field content into
the extracted modes.

\subsection{Diagnostics}
\label{sec:diagnostics}
Seven quantities are computed per frequency and stored with every
T-matrix:
\begin{center}
\begin{tabular}{lll}
\toprule
Quantity & Definition & Healthy value\\
\midrule
residual & $\|\mathsf{T}\mathsf{A}-\mathsf{F}\|_F / \|\mathsf{F}\|_F$ &
 solver accuracy ($10^{-3}$--$10^{-2}$)\\
conditioning & $\mathrm{cond}(\mathsf{A})$ & $<10^2$ (full Fibonacci:
 $\approx 1.3$; $C_{4v}$-reduced: $10$--$13$)\\
passivity & $\max \mathrm{sv}(\mathsf{S})$, $\mathsf{S} = \mathsf{I} +
 2\mathsf{T}$ & $\le 1 + \text{noise floor}$ ($\lesssim 1.03$
 observed)\\
unitarity & $\max|\mathrm{sv}(\mathsf{S}) - 1|$ & $\approx 0$ (lossless
 only)\\
reciprocity & $\|\mathsf{T} - \mathcal{R}(\mathsf{T})\|_F /
 \|\mathsf{T}\|_F$ & residual level near resonance$^{\dagger}$\\
incident deviation & measured vs.\ analytic $\mathbf{a}$ &
 $< 10^{-2}$\\
noise-order content & $\|\mathsf{T}\|_F$ fraction above the
 per-frequency & $\approx 0$ where $w(ka) \ge l_{\max}$\\
 & Wiscombe order $w(k a)$ & \\
\bottomrule
\end{tabular}
\end{center}
$^{\dagger}$The extraction error is an \emph{absolute} quantity, set by
the field error relative to the incident-field magnitude
(Sec.~\ref{sec:noisefloor}).  Far below resonance $\|\mathsf{T}\|$ is
small while this absolute error is unchanged, so the relative
quantities (reciprocity, residual) grow toward the low-frequency end
of the range even though the underlying data are no worse there.
Values at the ends of the range should be read against this
expectation.
The reciprocity map is
\begin{equation}
\mathcal{R}(\mathsf{T})_{(p',l',m'),(p,l,m)}
 = (-1)^{m+m'}\, \mathsf{T}_{(p,l,-m),(p',l',-m')},
\label{eq:reciprocity}
\end{equation}
which any scatterer composed of reciprocal media must satisfy to within
the extraction accuracy, independent of shape.  The sign structure of
Eq.~(\ref{eq:reciprocity}) was determined numerically by requiring
consistency with the amplitude reciprocity theorem
$\ehat_2\cdot\mathbf{F}(\khat_2\leftarrow\khat_1;\ehat_1) =
\ehat_1\cdot\mathbf{F}(-\khat_1\leftarrow-\khat_2;\ehat_2)$
\cite{Saxon1955} evaluated with the package's own plane-wave and
far-field expansions; alternative sign choices fail this test at order
unity.  A reciprocity violation at order unity therefore indicates an
extraction or convention error rather than physics.  Passivity and
unitarity follow from energy conservation \cite{Waterman1971}.

\subsection{Analytic references: homogeneous and coated spheres}
For validation, the diagonal T-matrix of a homogeneous sphere is computed
from the Mie coefficients $a_l, b_l$ of Bohren and Huffman
\cite{BohrenHuffman}, and that of a core--shell sphere from their
coated-sphere generalization.  The literature forms are written in the
$e^{-i\omega t}$ convention; in the package convention the entries are
\begin{equation}
T^{MM}_l = -\,b_l^{*}, \qquad T^{NN}_l = -\,a_l^{*},
\label{eq:mieconj}
\end{equation}
with complex refractive indices conjugated on input (loss corresponds to
a negative imaginary part in the $e^{+j\omega t}$ convention).  The
conjugation in Eq.~(\ref{eq:mieconj}) cannot be detected by internal
consistency tests, which are invariant under conjugation of the
reference; it was fixed by comparison with the full-wave solution of
the $\varepsilon=4$ sphere (Sec.~\ref{sec:validation}).  The
coated-sphere implementation is verified in its two exact limits (shell
index equal to the background: homogeneous sphere of the core; core
index equal to the shell: homogeneous sphere of the outer radius) and by
passivity for lossy shells.

\subsection{Conversion to the tmat.h5 convention}
\label{sec:treamsconv}
The data format \cite{Asadova2025format} and \emph{treams}
\cite{Beutel2024treams} use the physics convention $e^{-i\omega t}$ with
outgoing waves $h_l^{(1)}$, Jackson-normalized harmonics, and a basis
that differs from Sec.~\ref{sec:conventions} by an overall factor $i$,
which cancels in the T-matrix.  Writing one physical monochromatic field
in both phasor conventions, $\EE(\rr,t) = \mathrm{Re}[\EE_+ e^{+j\omega
t}] = \mathrm{Re}[\EE_- e^{-i\omega t}]$ with $\EE_- = \EE_+^*$, and
using $\Ptil_l^{-m} = (-1)^m \Ptil_l^m$, the basis functions conjugate as
$(\Mv{c}_{lm})^* = (-1)^m \Mv{c'}_{l,-m}$ (radial kinds mapping
$j_l \to j_l$, $h^{(2)}_l \to h^{(1)}_l$), so coefficients map
antilinearly, $a^-_{lm} = (-1)^m (a^+_{l,-m})^*$, and
\begin{equation}
\mathsf{T}^-_{(p,l,m),(p',l',m')}
 = (-1)^{m+m'} \left[\mathsf{T}^+_{(p,l,-m),(p',l',-m')}\right]^{*}.
\label{eq:convmap}
\end{equation}
Equation~(\ref{eq:convmap}) is \emph{not} element-wise conjugation; the
two coincide only for matrices that are diagonal and even in $m$, such
as spheres.  Applying the map twice returns the original matrix, and
the test suite verifies the map at the level of the fields themselves:
the complex conjugate of a field evaluated from arbitrary coefficients
equals the field evaluated from the mapped coefficients, and
transforming a T-matrix and transforming its input/output coefficient
vectors give the same scattered field.  Files on disk are written in
the converted convention with explicit mode tables; the reader applies
the inverse map, so quantities analyzed inside the package and
quantities read from disk agree.

%======================================================================
\section{Simulation setup in CST Studio Suite}
\label{sec:cstsetup}
The extraction imposes the following configuration, generated as history
blocks through the CST scripting interface.
\begin{description}
\item[Solver.]  The frequency-domain solver (FDFD) on a tetrahedral
mesh with second-order elements.  (CST's time-domain solver, an
FDTD-type method, is not used here.)  The active solver must be
switched explicitly (\texttt{ChangeSolverType}); a new project defaults
to the time-domain solver, and configuring the frequency-domain solver
alone does not change the active solver.
\item[Mesh strategy.]  \emph{Deterministic, non-adaptive} by default: a
fixed steps-per-wavelength density
(\texttt{ScattererPlan.mesh\_steps\_per\_wave}, with the
\texttt{CellsPerWavelengthPolicy} pinned --- see
Appendix~\ref{app:cstbugs}) and \texttt{mesh\_adaption=False}.  CST
re-runs adaptive refinement from scratch on \emph{every} solver start
and never carries an adapted mesh across solves
(Appendix~\ref{app:cstbugs}), so with adaptation enabled each
illumination pays the full multi-pass cost and each column of the
T-matrix is measured on a different mesh.  A fixed density makes every
solve identical and was validated on this pipeline against both an
adaptive reference and the exact Mie solution at matched accuracy and
roughly a quarter of the solve time.  The density must be
\emph{calibrated per geometry class} by a short convergence check
(compare the scattered column across densities and against one adaptive
run; \texttt{examples/density\_check\_spoke\_wheel.py} is the template).
\item[Splitting the frequency range.]  Over a wide spectral range the
monitor radius is fixed by the lowest frequency
(Sec.~\ref{sec:monitorplacement}) while the mesh density is fixed by
the highest, and a single project must accommodate both: a large
domain meshed finely.  For an extraction spanning 10--34~THz this
combination produced a mesh of order $10^7$ cells, and the solver ran
out of memory while constructing the preconditioner.  The remedy is to
divide the frequency range into sub-ranges, each with its own monitor
radius and mesh (\texttt{range\_edges\_thz} in the configuration;
\texttt{BAND\_EDGES\_THZ} in the extraction scripts); the same
10--34~THz extraction split at 20~THz required well under $10^6$ cells
per sub-range.  Each sub-range is an independent extraction writing
its own file; the files are merged afterwards by concatenation along
the frequency axis (\texttt{examples/merge\_band\_tmatrices.py}, which
verifies identical mode bases and non-overlapping frequency grids).
After merging, the continuity of $\mathsf{T}(f)$ at the junction of
two sub-ranges is evaluated automatically and reported as a
diagnostic: the two sides of the junction were computed with different
meshes and different monitor radii, so a step in $\mathsf{T}(f)$ there
that is no larger than the steps between neighboring frequencies
within a sub-range shows that the result is insensitive to both
choices.
\item[Excitation.]  Plane-wave source with linear polarization; the
propagation direction and polarization vector are set per illumination.
CST normalizes the excitation to the specified electric field vector
(V/m) with phase reference at the global origin.
\item[Boundaries and background.]  Open boundaries with added space set
\emph{explicitly on all six faces}, and the background material set
explicitly to vacuum.  Both are set explicitly rather than left to
defaults so that the configuration is complete and reproducible from
the history alone.  Note that in scripted use the ``apply in all
directions'' option does not copy a single face's added space to the
other faces (Appendix~\ref{app:cstbugs}); the resulting asymmetric
domain affected all early extractions until each face was set
individually.  No symmetry planes are used: the illumination set spans
arbitrary directions, and symmetry is exploited algebraically instead
(Sec.~\ref{sec:symmetryreduction}).
\item[Monitors.]  One electric-field and one magnetic-field monitor per
extraction frequency, with the monitor frequencies forced as exact solver
samples (never interpolated).  For plane-wave excitation the results
appear in the navigation tree with the label suffix \texttt{[pw]}.
\item[Field export.]  Complex E and H are exported at the surface
quadrature nodes through the \texttt{ASCIIExport} object with a
user-supplied evaluation-point file.  Exported coordinates are compared
against the requested points, with automatic detection of a global unit
scale; a mismatch aborts the run.  Magnetic fields are exported in A/m
and multiplied by $Z_0$.
\item[Result lifetime and caching.]  Redefining the plane wave
invalidates solver results, so the fields of each illumination are
exported immediately after its solve.  Completed illuminations are
written to disk (\texttt{.npz}), and an interrupted extraction resumes
from them without recomputation.  Each cached record stores the
frequency grid, the illumination direction and polarization, and the
monitor radius, and is reused only when all of these match the current
run.  This matters because the illumination list and the monitor
radius can legitimately differ between two runs of the same object ---
symmetry reduction changes the list of solved directions, and a
revised monitor radius changes the surface the fields were recorded
on --- so a cache matched by position in the list alone could supply a
column belonging to a different direction or a different monitor.
\item[Run log.]  Everything printed during an extraction (including
warnings and the diagnostics table) is appended to
\texttt{exports/log.txt} in the run directory with a timestamped header
per invocation.
\end{description}

%======================================================================
\section{Software}

\subsection{Modules}
\begin{center}
\begin{tabular}{ll}
\toprule
Module & Content\\
\midrule
\texttt{vswf.py} & basis functions, projections, plane-wave expansion,\\
 & separation, convention maps\\
\texttt{quadrature.py} & surface quadrature; Fibonacci and Lebedev
 direction sets\\
\texttt{mie.py} & homogeneous and coated sphere references;
 Wiscombe criterion\\
\texttt{tmatrix\_solve.py} & least-squares solve; passivity, unitarity,
 reciprocity\\
\texttt{storage.py} & tmat.h5 writing and reading with convention
 conversion\\
\texttt{pipeline.py} & orchestration: project setup, illumination loop,
 caching,\\
 & monitor-placement criteria, noise-order diagnostics\\
\texttt{config.py} & user settings (directories), extraction options,
 automatic-check thresholds\\
\texttt{cst\_driver/} & CST session control, geometry builders, template
 mode,\\
 & excitation, monitors, solver configuration, field export\\
\texttt{postprocess/symmetry.py} & point-group representations,
 symmetry-reduced\\
 & illumination, STL point-group detector, group-average denoising\\
\texttt{postprocess/metasurface.py} & forward-scattering amplitude and
 specular\\
 & S-parameters (isolated-scatterer approximation)\\
\texttt{postprocess/coupling.py} & finite-cluster Foldy--Lax:
 translation operators,\\
 & cluster solve, separation guards\\
\bottomrule
\end{tabular}
\end{center}
The package is self-contained: it depends only on \texttt{numpy},
\texttt{scipy}, \texttt{h5py}, and the CST Python interface.

\subsection{Installation and configuration}
\label{sec:install}
Create a Python environment ($\ge$ 3.9) with \texttt{numpy},
\texttt{scipy}, \texttt{h5py}, and \texttt{pytest}, and install the CST
interface shipped with the CST installation:
\begin{verbatim}
pip install --no-index --find-links \
    "<CST folder>/Library/Python/repo/simple" cst-studio-suite-link
\end{verbatim}
Machine-specific settings are collected in an optional JSON file
\texttt{settings.json} at the repository root, written once per
machine:
\begin{verbatim}
{
  "run_root": "D:/cst_runs",
  "library": "",
  "max_cpus": 6,
  "hardware_acceleration": false
}
\end{verbatim}
\texttt{run\_root} (CST projects and solver output) must be a local
disk, because cloud-synchronized folders lock files during solves;
\texttt{library} (destination of the final \texttt{.tmat.h5} files)
defaults to \texttt{library/} in the repository.  The CPU count and
the hardware-acceleration setting belong here because they differ
between machines while the physics configuration does not; an
extraction configuration can override them for a single run.  The
environment variables \texttt{CST\_TMATRIX\_RUN\_ROOT} and
\texttt{CST\_TMATRIX\_LIBRARY} take precedence over
\texttt{settings.json} where both are set.

\subsection{Operating procedure}
\begin{enumerate}
\item \textbf{Offline verification} (no CST required):
\texttt{python -m pytest tests/ -q}.  All tests must pass before any
extraction.
\item \textbf{Single-solve verification} (requires CST):
\texttt{python examples/smoke\_test\_sphere.py} builds a dielectric
sphere, runs one solve, and checks the measured incident wave against the
analytic plane-wave expansion and the measured scattered wave against Mie
theory.  Run this once per machine or CST version.
\item \textbf{Extraction.}  Describe the scatterer by a
\texttt{ScattererPlan} and call \texttt{extract\_tmatrix} (see
Sec.~\ref{sec:usage}).
\end{enumerate}

\subsection{Automatic checks and unattended operation}
\label{sec:rungates}
An extraction takes hours, but most of the conditions that invalidate
one can be tested in seconds to minutes.  \texttt{extract\_tmatrix}
therefore performs three checks automatically (thresholds in
\texttt{ExtractionConfig}; a failed check raises
\texttt{ExtractionCheckError} and stops the run):
\begin{description}
\item[Before any computation.]  The two monitor-placement conditions of
Sec.~\ref{sec:monitorplacement} are required, not merely reported.  The
automatically computed monitor factor satisfies them by construction;
the check protects against configurations that bypass the automatic
factor --- an explicit \texttt{monitor\_factor} in a script, an old
configuration file, or a future code change --- reaching the solver.
\item[After meshing, before any solve.]  The mesh cell count and the
estimated solver memory ($\approx$12~kB per second-order tetrahedron;
Appendix~\ref{app:cstbugs}) are compared with configured limits and
with the memory currently available.  A configuration that would
exhaust memory is rejected after the few minutes of meshing rather
than discovered when the solver fails.
\item[After every solve.]  The incident deviation of the completed
solve is evaluated immediately.  Above $10^{-2}$ (default) it is
logged; above $3\times10^{-2}$ (default) the extraction stops, since a
recorded field this far from the known incident wave invalidates the
column regardless of the cause --- including causes not yet
understood.  Columns from a stopped solve are not written to the
cache, so a corrected rerun cannot reuse them unnoticed.  After the
first solve the measured top-order content
(Sec.~\ref{sec:truncation}) and an estimate of the total run time are
also printed, the latter with an optional upper limit.
\end{description}
As a point of reference: the per-solve check detected the misplaced
monitor of Sec.~\ref{sec:monitorplacement} after 21 minutes; the same
physical mistake, made before these checks existed, was found only
after a completed 5.5-hour run.

For extracting several designs without supervision,
\texttt{examples/batch\_extract\_spoke\_wheel.py} runs the designs
sequentially, verifies the placement conditions for every sub-range
before CST is started, opens a fresh CST connection per sub-range (so
a crash of the CST front end in one does not affect the others),
continues with the remaining work if one fails, merges each design's
sub-ranges, evaluates the continuity of $\mathsf{T}(f)$ at their
junctions, and appends a summary to \texttt{batch\_log.txt}.  A
sub-range stopped by the incident-deviation check is rerun once
automatically with the monitor radius increased by one Wiscombe order;
a second stop is left for the operator.

Solver resources (thread count, hardware acceleration) are properties
of the machine, not of the physics problem, and are therefore read
from the per-machine \texttt{settings.json} (Sec.~\ref{sec:install})
rather than fixed in the extraction configuration; a configuration may
still override them for a single run.

%======================================================================
\section{Usage}
\label{sec:usage}

\subsection{Scripted geometries}
\begin{verbatim}
import numpy as np
from cst_tmatrix import ScattererPlan, ExtractionConfig, extract_tmatrix
from cst_tmatrix.cst_driver import CSTSession, builder

def build(h):
    builder.define_material(h, "si", eps=11.7)
    builder.new_component(h, "scatterer")
    builder.build_cylinder(h, radius=30.0, zmin=-25.0, zmax=25.0,
                           material="si")     # project length units

plan = ScattererPlan(name="si_disk", build=build, r_circ_m=39.1e-6,
                     freqs_hz=np.array([0.6e12, 0.8e12, 1.0e12]))
T, diagnostics, path = extract_tmatrix(CSTSession(), plan)
\end{verbatim}
Arbitrary geometry is supported through
\texttt{h.add\_history(caption, vba)} with any valid CST history block.
The geometry must be centered at the origin, which is the origin of the
multipole expansion.

\subsection{Template projects}
\label{sec:template}
Complex geometries are often easier to build interactively.  In template
mode the user supplies a finished CST project containing \emph{only} the
geometry and materials:
\begin{verbatim}
plan = ScattererPlan(name="my_object", template=r"path\to\object.cst",
                     r_circ_m=50e-6, freqs_hz=np.array([1.0e12]),
                     length_unit="um", freq_unit="THz")
\end{verbatim}
The file is copied into the run directory (the original is never
modified), verified, and completed with the extraction settings of
Sec.~\ref{sec:cstsetup}.  Requirements on the template:
\begin{itemize}
\item geometry centered at the origin;
\item no ports and no periodic or unit-cell boundaries;
\item no structure extending to the domain boundary (in particular no
substrate slab): the monitor sphere must lie in homogeneous background;
\item built in the declared length and frequency units (units are not
modified, because reinterpreting units would rescale the geometry).
\end{itemize}
The verification reports PASS/FAIL/UNVERIFIED per check; the bounding box
of all solids is compared against the declared circumscribing radius,
which catches decentered geometry and unit or radius mistakes.  Checks
that the scripting interface cannot perform are reported as UNVERIFIED
for manual confirmation.

\subsection{Extraction from a configuration file}
\label{sec:noconfig}
An extraction can be specified entirely by a JSON configuration file;
no Python is written.  The geometry is built interactively in CST
(template requirements, Sec.~\ref{sec:template}), the file is filled
in following the annotated example
\texttt{examples/example\_extraction.json}, and the extraction is run
with
\begin{verbatim}
python -m cst_tmatrix.runner  my_extraction.json --dry-run
python -m cst_tmatrix.runner  my_extraction.json
\end{verbatim}
With \texttt{--dry-run} the configuration is validated and the
monitor-placement quantities are printed for every frequency
sub-range; CST is not started.  Without it, the sub-ranges are
extracted in sequence with the defaults of this manual (automatic
monitor placement, fixed mesh density, the automatic checks, symmetry
reduction if configured), merged, tested for continuity at the
sub-range junctions, and written as one \texttt{.tmat.h5} file.

The configuration file contains the quantities that describe the
problem: the geometry file, the circumscribing radius, the frequency
grid and its subdivision, $l_{\max}$, and optionally the point-group
symmetry of the cell.  Machine-dependent settings (working directory,
CPU count, hardware acceleration) are kept separately in
\texttt{settings.json}, written once per machine
(Sec.~\ref{sec:install}).  Since JSON has no comment syntax, keys
beginning with an underscore are ignored and may hold notes; the
example file uses them for its documentation.  An invalid or
incomplete configuration is reported at load time with a message
naming the offending entry.

\subsection{Extraction options}
\texttt{ExtractionConfig} sets the truncation order (\texttt{lmax}),
the overdetermination factor (\texttt{illumination\_factor}), the
distribution of incidence directions (\texttt{illumination\_scheme} =
\texttt{"fibonacci"} or \texttt{"lebedev"}), the surface-quadrature
oversampling, the symmetry reduction
(Sec.~\ref{sec:symmetryreduction}), and the thresholds of the
automatic checks: the mesh-size and memory limits, the two
incident-deviation levels, and an optional limit on the projected run
time.  The monitor factor is computed automatically by default
(Sec.~\ref{sec:monitorplacement}); the mesh is the fixed density of
Sec.~\ref{sec:cstsetup}; a wide frequency range should be subdivided
as described there.

Before the first long extraction of a new class of geometry, two short
preparatory computations are worthwhile.  A mesh-density comparison
(solve one illumination at two or three densities and against one
adaptive run, and compare the scattered coefficients;
\texttt{examples/density\_check\_spoke\_wheel.py} shows the procedure)
establishes the density the geometry requires.  A timing run of two
illuminations at the full configuration then gives the incident
deviation and the time per solver run, from which the duration of the
complete extraction follows by multiplication.

\subsection{Output files}
Each extraction writes one HDF5 file in the tmat.h5 layout:
\begin{center}
\begin{tabular}{ll}
\toprule
Element & Content\\
\midrule
\texttt{/tmatrix} & $(n_f, N, N)$ complex, convention of
 Sec.~\ref{sec:treamsconv}\\
\texttt{/frequency} & frequencies with \texttt{unit} attribute (Hz)\\
\texttt{/modes/l, m, polarization} & mode tables;
 \texttt{"electric"}/\texttt{"magnetic"} parity labels\\
\texttt{/embedding} & background medium (vacuum)\\
\texttt{/scatterer} & geometry and material metadata\\
\texttt{/computation} & method, parameters, and
 \texttt{analysis/} diagnostics\\
root attributes & name, description, keywords, format version\\
\bottomrule
\end{tabular}
\end{center}
The mode tables make the file self-describing; the reader
(\texttt{storage.load\_tmatrix}) accepts any parity-basis ordering and
reconstructs the package-internal representation.

%======================================================================
\section{Validation}
\label{sec:validation}

\subsection{Offline tests}
The suite (92 tests, no CST required) verifies: normalized Legendre and
angular functions against \texttt{scipy}; ladder identities against
finite differences; projection round trips at $10^{-13}$; the plane-wave
expansion at $10^{-13}$; synthetic T-matrix recovery from fields with
randomized source phases at $4\times 10^{-15}$; noise averaging
($10^{-3}$ field noise $\to$ $4\times 10^{-5}$ matrix error); the far-field
projection route; the reciprocity map against the amplitude reciprocity
theorem at $10^{-12}$; the coated-sphere reference in both exact limits;
passivity and unitarity of lossy and lossless references; Lebedev rule
exactness; the convention conversion at the field level; storage round
trips through HDF5; the monitor-placement criteria against the recorded
recorded placement observations (Sec.~\ref{sec:monitorplacement}); the
point-group representations, orbit conventions, orbit degeneracy, and
conditioning of the reduced illumination system
(Sec.~\ref{sec:symmetryreduction}), including full-rank recovery of a
random $C_{4v}$-symmetric matrix from the reduced set at $10^{-8}$; the
noise-order diagnostics; and the finite-cluster Foldy--Lax machinery
(translation-operator field reproduction, $\rho$-independence, Born
limit, and separation guards, Sec.~\ref{sec:metasurfacepost}).

One limitation of synthetic testing should be stated explicitly:
noise-free recovery tests do not probe conditioning.  The clustered
illumination directions of Sec.~\ref{sec:symmetryreduction} reproduced
synthetic T-matrices to $10^{-8}$, yet the same construction failed on
measured data with a field error of only $10^{-3}$, because the column
system had condition number $4\times10^{5}$.  The suite therefore
asserts condition numbers in addition to recovery errors.

\subsection{Comparison with full-wave reference solutions}
A dielectric sphere ($\varepsilon = 4$, radius 40~\textmu m,
$ka = 0.63$--$1.26$) was extracted with CST Studio Suite 2026 at 0.5,
0.75, and 1.0~THz using $l_{\max}=3$ (46 solver runs).  The comparison
with the analytic Mie T-matrix gives:
\begin{center}
\begin{tabular}{cccccc}
\toprule
$f$ (THz) & residual & cond($\mathsf{A}$) & $\max\mathrm{sv}(\mathsf{S})$
 & incident dev. & max.\ rel.\ error vs.\ Mie\\
\midrule
0.50 & $1.3\times10^{-3}$ & 1.3 & 1.000066 & $3.8\times10^{-2}$ &
 $3.4\times10^{-3}$\\
0.75 & $1.0\times10^{-3}$ & 1.3 & 1.000174 & $1.4\times10^{-2}$ &
 $3.3\times10^{-3}$\\
1.00 & $9.2\times10^{-4}$ & 1.3 & 1.000385 & $9.8\times10^{-3}$ &
 $3.0\times10^{-3}$\\
\bottomrule
\end{tabular}
\end{center}
Off-diagonal elements, which vanish analytically for a sphere, are
recovered at $3$--$4$ orders of magnitude below the dominant elements;
they measure the extraction noise floor, since the least-squares solve
imposes no symmetry.  The single-solve verification recovers the incident
wave with amplitude $1.0013$ and phase $6\times10^{-4}$~rad, confirming
CST's documented source normalization, and matches the Mie prediction of
the scattered coefficients to $1.6\times10^{-3}$.

This comparison exposed two convention errors that offline testing could
not: the missing conjugation of the Mie reference,
Eq.~(\ref{eq:mieconj}), and the necessity of switching the active solver
explicitly (Sec.~\ref{sec:cstsetup}).  Both are corrected and documented;
the general lesson is that internal consistency tests are blind to any
error that amounts to a conjugation of the reference, and at least one
comparison against an independent full-wave solution is required to fix
the convention.

\subsection{The absolute scale of the extraction error}
\label{sec:noisefloor}
The error of the extracted coefficients is set by an absolute scale,
not a relative one.  The recorded total field is dominated by the
incident wave, and the solver's relative error $\epsilon \approx
10^{-3}$ on that field enters the projected incident and scattered
coefficients with the same \emph{absolute} magnitude,
$\sim\epsilon\,\|\mathbf{a}\|$.  This was verified on a weakly
scattering sphere ($|T| \approx 0.025$): at every multipole order, the
absolute error of the measured incident coefficients (whose exact
values are known) and of the scattered coefficients agreed within a
factor of two.

Two consequences follow for the interpretation of comparisons.
First, for a weak scatterer the ratio of incident to scattered
amplitude is large ($255{:}1$ in the case above), so the \emph{relative}
error of the scattered coefficients is large even when the extraction
is working correctly: an apparent $19\%$ disagreement with the exact
Mie solution consisted of dominant coefficients agreeing to
$0.5$--$3\%$ together with coefficients whose exact values lie below
the absolute error scale and which therefore contain only noise.
Second, the same effect appears at the off-resonant ends of a
resonator's frequency range, where $\|\mathsf{T}\|$ is small
(Sec.~\ref{sec:diagnostics}).  A single relative norm over an
oversized mode space is therefore not a meaningful error measure;
comparisons should be made order by order against the measured
absolute scale.

\subsection{Benchmark against an independent solver: TiO$_2$ cylinder}
\label{sec:cylinderbench}
The end-to-end pipeline was validated against the \emph{treams}
example T-matrix of a TiO$_2$ cylinder (radius 250~nm, height 300~nm,
$n = 2.5$, vacuum), computed independently with JCMsuite
(rotationally-symmetric 2D FEM with built-in multipole expansion) ---
a different solver, mesh type, and extraction method, at a strong
scatterer ($\|\mathsf{T}\| \approx 2.5$--$2.9$), for which the absolute
error scale of Sec.~\ref{sec:noisefloor} is far below the signal.
Extraction: the CST frequency-domain solver on a tetrahedral mesh at
$l_{\max} = 4$ (matching the reference), $C_{4v}$
symmetry reduction (10 solver runs instead of 72), deterministic
non-adaptive mesh, at 320~THz (off-resonance) and 393.6~THz (the
strongest extinction resonance in the range).  Two independent comparisons:
comparison:
\begin{enumerate}
\item \emph{Column-wise, no fit involved}: every solved illumination
satisfies $\mathbf{f}_{\mathrm{meas}} =
\mathsf{T}_{\mathrm{ref}}\,\mathbf{a}_{\mathrm{meas}}$ to
$0.4$--$0.8\%$ --- a direct test of the CST solve, field export,
projection, and separation chain against independent physics, with no
free parameters.
\item \emph{Matrix-wise}: relative Frobenius difference $0.73\%$
(320~THz) and $1.13\%$ (393.6~THz); the six dominant entries agree to
$0.1$--$0.2\%$ individually; the extinction observable to
$0.15$--$0.5\%$.  Two structural properties \emph{not} imposed by the
extraction emerged from the data: $m$-conservation (the $C_{4v}$
expansion only enforces $m - m' \equiv 0 \bmod 4$; the continuous
rotational symmetry appeared at the $0.5$--$0.8\%$ level) and
reciprocity ($0.9$--$1.4\%$).  Entries that are structural zeros in the
2D reference are recovered at a median of $2\times 10^{-16}$
($C_{4v}$-forbidden, enforced) or below $8.5\times 10^{-3}$
(allowed-but-unpopulated, the measured noise floor).
\end{enumerate}

\subsection{Recommended validation sequence for new object classes}
Each step adds one physical feature with an available check: a lossy
sphere (exact Mie; passivity strictly below unity); a coated sphere (the
last object with an exact reference); a finite cylinder --- \emph{done},
Sec.~\ref{sec:cylinderbench}, against an independent full-wave reference
rather than only a selection rule; a sphere dimer (comparison with the
translation-addition composition of the single-sphere T-matrix, now
available offline in \texttt{postprocess/coupling.py}); a cube
($C_{4v}$ selection rules; mesh convergence at edges); a helix
(reciprocity nontrivial; magneto-electric coupling blocks).  For
extractions performed in sub-ranges, the junction continuity test of
Sec.~\ref{sec:cstsetup} applies to every new design at no cost.

%======================================================================
\section{Sources of error and how they manifest}
\label{sec:failuremodes}
Each mechanism below has been observed on this pipeline.  Several of
them produce superficially similar diagnostics; the frequency
dependence and the multipole distribution of the anomaly usually
identify the cause.  The diagnostics of a stored matrix should be read
against this list before the matrix is used.

\begin{description}
\item[Monitor inside the reactive near field (condition ii,
Sec.~\ref{sec:monitorplacement}).]  Incident deviation of order
$10^{-1}$, roughly constant across the frequency range and largest at the
\emph{low}-frequency end; singular values of $\mathsf{S}$ well above
1; the $l = l_{\max}$ block of $\mathsf{T}$ an order of magnitude too
large; and, in a merged file, a jump of an order of magnitude in
reciprocity at the junction of two sub-ranges.  The junction
discontinuity and
the low-frequency weighting distinguish this from insufficient mesh
resolution.
\item[Ill-conditioned decomposition at $l = l_{\max}$ (condition i).]
Incident deviation up to order unity, concentrated at $l = l_{\max}$
and worst at the lowest frequency; occurs when $k_{\min}R$ is at or
below $l_{\max}$.
\item[Undetermined multipole orders.]  Per-order content that
\emph{grows} with $l$ and is nearly frequency-independent.  The
physical orders are unaffected --- each $(l,m)$ pair is determined by
an independent $2\times 2$ system --- but norms and sums over all
modes are contaminated.  Occurs when $l_{\max}$ is set well above the
measurable content (Sec.~\ref{sec:truncation}).
\item[Small-$\|\mathsf{T}\|$ inflation of relative quantities.]
Residual and reciprocity growing toward the off-resonant ends of the
range while
the incident deviation does not; a consequence of the absolute error
scale, not a defect (Sec.~\ref{sec:noisefloor}).
\item[Insufficient mesh resolution at the upper end of the range.]  All
diagnostics degrading smoothly and monotonically toward the highest
frequency, most strongly for the electrically largest object of a
family; observed up to incident deviation $1.6\times10^{-2}$ and
$\max\mathrm{sv}(\mathsf{S}) = 1.05$ at an upper range limit $2.7\times$ above
the design resonance.  The direction of the trend separates this from
the misplaced monitor: mesh resolution is worst at the top of the
range and varies smoothly; the misplaced monitor is worst at the
bottom,
nearly flat, an order of magnitude larger, and discontinuous at a
sub-range junction.  If the upper end of the range matters for the
application, raise the mesh density for that sub-range; otherwise the stored
diagnostics record the magnitude of the effect.
\item[Ill-conditioned reduced illumination set.]  Plane-wave
observables (e.g.\ extinction) correct while the matrix violates
reciprocity, passivity, and the object's known selection rules;
condition number of the expanded column system $\gg 10^2$ with
numerical rank below $N$.  Occurs when the solved directions cluster
on the sphere (Sec.~\ref{sec:symmetryreduction}).
\item[Insufficient memory (``Could not compute preconditioner'').]
Meshing succeeds and the first frequency sample fails.  Produced by
the combination of a large monitor radius (set by the lowest
frequency) with a fine mesh (set by the highest) over a wide spectral
range in one project.  Subdivide the range (Sec.~\ref{sec:cstsetup}); the estimate of
$\approx$12~kB per second-order tetrahedron against available memory
predicts the problem before solving.
\item[Columns computed on different meshes.]  Occurs if adaptive
refinement is left on (each solve adapts independently) or mesh
settings change mid-run.  There is no single sharp symptom, which is
why the mesh must be fixed before the illumination loop begins
(Sec.~\ref{sec:cstsetup}).
\item[Reuse of cached columns from a different configuration.]
Excluded by construction: cached columns are reused only when the
frequency grid, illumination direction, polarization, and monitor
radius all match (Sec.~\ref{sec:cstsetup}).
\end{description}

%======================================================================
\section{Planned extensions}
\label{sec:planned}

\subsection{Metasurface post-processing: implemented stages}
\label{sec:metasurfacepost}
Post-processing consumes stored T-matrices; the extraction machinery is
not involved.  Two stages are implemented and validated offline:

\emph{Specular S-parameters in the isolated-scatterer approximation}
(\texttt{postprocess/metasurface.py}).  With the package scattering
amplitude $\mathbf{F}$ defined by $\EE_{\mathrm{sca}} \to
\mathbf{F}(\theta,\varphi)\, e^{-jkr}/(kr)$, the forward-scattering
theorem gives, for a periodic array of cell area $A$ under incidence
$(\kk_{\mathrm{inc}}, \ehat_{\mathrm{inc}})$,
\begin{equation}
S_{21} = 1 - \frac{j\,2\pi}{k^{2}A\cos\theta_{\mathrm{inc}}}\,
 \ehat^{*}_{\mathrm{inc}}\cdot\mathbf{F}(\kk_{\mathrm{inc}}),
\qquad
S_{11} = -\frac{j\,2\pi}{k^{2}A\cos\theta_{\mathrm{inc}}}\,
 \ehat^{*}_{\mathrm{ref}}\cdot\mathbf{F}(\kk_{\mathrm{ref}}).
\label{eq:sparams}
\end{equation}
Both the $1/k^2$ (from the dimensionless normalization of $\mathbf{F}$
above) and the sign of $j$ (from the $e^{+j\omega t}$ convention) differ
from the textbook $e^{-i\omega t}$ form; each was initially transcribed
wrongly and fixed by an independent check --- the magnitude against a
direct CST periodic simulation, the sign by the optical theorem
$\mathrm{Im}\,[\ehat^*\!\cdot\mathbf{F}_{\mathrm{fwd}}]/k =
-(k/4\pi)\,\sigma_{\mathrm{ext}}$ evaluated on an exact Mie sphere.
Transcribing phasor-convention-sensitive formulas without such a check
is not safe in this codebase.  The isolated-scatterer approximation
itself carries a real physical error (inter-cell coupling is absent);
comparison against a direct periodic CST solve of a resonant cell shows
qualitative agreement with quantitative gaps that the Foldy--Lax stage
is designed to close.

\emph{Finite-cluster Foldy--Lax} (\texttt{postprocess/coupling.py}).
Self-consistent outgoing coefficients $\mathbf{b}_i =
\mathsf{T}_i\big(\mathbf{a}_i + \sum_{j\ne i}
\mathsf{C}(\rr_i - \rr_j)\,\mathbf{b}_j\big)$ for a finite cluster
\cite{Foldy1945,Tsang2001}.  The translation-addition operators
$\mathsf{C}$ are constructed \emph{numerically} from the package's own
validated field primitives (evaluate an outgoing mode about one centre,
re-project onto regular modes about the other) rather than from
closed-form Gaunt/Wigner coefficients --- the same
convention-risk argument as above.  Two measured facts govern accuracy:
the re-projection sphere must have the radius of the \emph{receiving
scatterer}, never a fraction of the separation (a regular expansion on
a sphere of radius $\rho$ carries angular content $\sim k\rho$, so
$\rho$ tied to the separation destroys the truncated expansion ---
measured error $0.66$ at $k\rho = 20$ with $l_{\max} = 8$ versus
$10^{-5}$--$10^{-6}$ with $\rho = a$); and accuracy degrades as
circumscribing spheres approach ($\sim$$10^{-3}$ at twice the touching
distance, $\sim$$6\times 10^{-2}$ at $1.3\times$), guarded by
\texttt{check\_pair\_separation}.

\subsection{Planned: periodic lattice sum and substrate coupling}
The infinite-array specular response requires the lattice sum
$\Omega = \sum_{\mathbf{R}\neq 0} \mathsf{C}(\mathbf{R})\,
e^{j\kk_\parallel\cdot\mathbf{R}}$, which is conditionally convergent
and requires Ewald summation; the per-site algebra
$\mathsf{T}_{\mathrm{eff}} = (\mathsf{I} -
\mathsf{T}\,\Omega)^{-1}\mathsf{T}$ is already implemented and tested
against its defining equation.  For a scatterer above a substrate, the
free-space T-matrix is coupled to the interface through a reflection
operator $\mathsf{R}$ constructed by expanding outgoing spherical waves
into plane waves (Weyl identity), applying Fresnel reflection, and
projecting back: $\mathsf{T}_{\mathrm{coupled}} = (\mathsf{I} -
\mathsf{T}\mathsf{R})^{-1}\mathsf{T}$.  The acceptance test for the
combined stage is the comparison of the composed unit-cell response
against a direct CST simulation of the same cell with periodic
(unit-cell) boundaries; CST's multilayered background for plane-wave
excitation offers an independent check of the substrate operator
alone.

%======================================================================
\appendix
\section{CST-specific behaviors and workarounds (bug report)}
\label{app:cstbugs}
Every item below was established by controlled experiment against CST
Studio Suite 2026 on this pipeline (dates given); none is documented in
the CST help in the form stated here.  They are collected in one place
because each silently corrupted or blocked real extractions before it
was found.

\begin{enumerate}
\item \textbf{\texttt{Background.ApplyInAllDirections} does not
broadcast in scripted history blocks} (2026-08-05).  Setting one face
distance (e.g.\ \texttt{.XminSpace}) with
\texttt{.ApplyInAllDirections "True"} pads \emph{only that face}; the
flag is a GUI convenience.  Verified via
\texttt{Boundary.GetCalculationBox()}: the domain was asymmetric
(zero-padded in $y$ and $z$) in every extraction until all six
\texttt{.XminSpace}\,\ldots\texttt{.ZmaxSpace} were emitted explicitly.
Workaround: \texttt{builder.set\_background\_vacuum} writes all six.

\item \textbf{\texttt{MeshSettings} steps-per-wavelength values are
silently ignored under the default policy} (2026-08-05).
\texttt{CellsPerWavelengthPolicy} defaults to \texttt{"automatic"}, in
which CST computes its own density and disregards
\texttt{StepsPerWaveNear/Far} entirely: settings of 15/6, 30/12, and
50/20 all produced the identical cell count until the policy was pinned
to \texttt{"cellsperwavelength"}, after which the same three settings
gave 82k, 599k, and 2.33M cells.  Any workflow that believes it is
controlling tetrahedral density without pinning this policy is not.

\item \textbf{No adaptive-mesh reuse across solver runs} (2026-08-05).
\texttt{FDSolver.MeshAdaptionTet} is re-evaluated on every
\texttt{.Start}; an adapted mesh is never carried to the next solve.
Three mechanisms were tested and refuted in continuous sessions with
the same project handle: (a) toggling the flag off after a converged
adaptive solve \emph{reverts to the coarse base mesh} (1{,}022{,}285
$\to$ 81{,}734 cells) rather than freezing the adapted one; (b)
changing the excitation via
\texttt{StoreDoubleParameter}+\texttt{RebuildOnParametricChange} (the
result-preserving parametric path) behaves identically; (c) driving the
sweep through CST's own \texttt{ParameterSweep} engine costs the full
per-point adaptation ($2\times$ single-solve time for 2 points).
Consequence: adaptation, if enabled, is paid per illumination on a
per-column mesh; the pipeline's default is the deterministic fixed
density of Sec.~\ref{sec:cstsetup} instead.

\item \textbf{\texttt{MeshAdaption3D.MinPasses} rejects 1}
(2026-08-05).  Values below 2 raise ``Please specify an integer number
greater or equal 2.''

\item \textbf{\texttt{MaxCPUs} requires \texttt{LimitCPUs "True"}}
(2026-08-03).  Without the (undocumented) \texttt{LimitCPUs} flag, the
thread count in \texttt{MaxCPUs} is ignored and all cores are used.
GPU offload uses the single-argument history form
\texttt{.HardwareAcceleration "True"}.

\item \textbf{Killing the controlling Python process does not abort an
in-flight solve} (2026-08-06).  The solver
(\texttt{Solver\_HF\_Tet\_FD\_AMD64}) is a separate process that
continues to completion, holding CPU cores and memory; two live
incidents contaminated subsequent timing runs.  After stopping a run
mid-solve, check for and terminate the solver process explicitly.

\item \textbf{``Could not compute preconditioner'' is the
out-of-memory signature} (2026-08-06).  Meshing succeeds, the first
frequency sample dies.  Observed calibration: $\sim$12~kB per
second-order tetrahedron; a $\sim$$10^7$-cell configuration exceeds a
32~GB machine.  See the subdivision of the frequency range,
Sec.~\ref{sec:cstsetup}.

\item \textbf{\texttt{PlaneWave.Normal}/\texttt{.EVector} accept
project-parameter names despite double-typed signatures} (2026-08-05).
The general rule that VBA \texttt{double} parameters reject parametric
expressions does not hold here; \texttt{GetNormal()} confirms the
values track \texttt{StoreDoubleParameter} updates after
\texttt{RebuildOnParametricChange}.  (Exploited only diagnostically;
the pipeline redefines the plane wave per illumination.)

\item \textbf{\texttt{Mesh.Update()} silently returns zero cells
without the high-frequency solver context} (2026-08-06).  Callable
only after \texttt{ChangeSolverType} and solver configuration;
\texttt{ForceUpdate()} alone invalidates without rebuilding.  Also,
\texttt{Mesh.Update} does \emph{not} pick up changed
\texttt{MeshSettings} densities in an already-meshed session (item~2
interacts); a fresh project reflects them.

\item \textbf{Operational: Python stdout block-buffering hides live
progress} (2026-08-05).  When the interpreter's output is piped (logs,
IDE consoles), prints buffer for minutes-to-hours; run extraction
scripts with \texttt{python -u}.
\end{enumerate}

%======================================================================
\begin{thebibliography}{99}
\bibitem{Waterman1965} P.~C.~Waterman, ``Matrix formulation of
electromagnetic scattering,'' Proc.\ IEEE \textbf{53}, 805--812 (1965).
\bibitem{Waterman1971} P.~C.~Waterman, ``Symmetry, unitarity, and
geometry in electromagnetic scattering,'' Phys.\ Rev.\ D \textbf{3},
825 (1971).
\bibitem{Mishchenko2002} M.~I.~Mishchenko, L.~D.~Travis, and
A.~A.~Lacis, \emph{Scattering, Absorption, and Emission of Light by
Small Particles} (Cambridge University Press, 2002).
\bibitem{Foldy1945} L.~L.~Foldy, ``The multiple scattering of waves,''
Phys.\ Rev.\ \textbf{67}, 107--119 (1945).
\bibitem{Tsang2001} L.~Tsang, J.~A.~Kong, and K.-H.~Ding,
\emph{Scattering of Electromagnetic Waves} (Wiley, 2001).
\bibitem{Fruhnert2017} M.~Fruhnert, I.~Fernandez-Corbaton,
V.~Yannopapas, and C.~Rockstuhl, ``Computing the T-matrix of a
scattering object with multiple plane wave illuminations,'' Beilstein
J.\ Nanotechnol.\ \textbf{8}, 614--626 (2017).
\bibitem{Jackson1999} J.~D.~Jackson, \emph{Classical Electrodynamics},
3rd ed.\ (Wiley, 1999).
\bibitem{BohrenHuffman} C.~F.~Bohren and D.~R.~Huffman, \emph{Absorption
and Scattering of Light by Small Particles} (Wiley, 1983).
\bibitem{Wiscombe1980} W.~J.~Wiscombe, ``Improved Mie scattering
algorithms,'' Appl.\ Opt.\ \textbf{19}, 1505--1509 (1980).
\bibitem{Lebedev1976} V.~I.~Lebedev, ``Quadratures on a sphere,''
USSR Comput.\ Math.\ Math.\ Phys.\ \textbf{16}, 10--24 (1976).
\bibitem{Saxon1955} D.~S.~Saxon, ``Tensor scattering matrix for the
electromagnetic field,'' Phys.\ Rev.\ \textbf{100}, 1771 (1955).
\bibitem{Asadova2025format} N.~Asadova \emph{et al.}, ``T-matrix
representation of optical scattering response: Suggestion for a data
format,'' J.\ Quant.\ Spectrosc.\ Radiat.\ Transf.\ \textbf{333},
109310 (2025).
\bibitem{Beutel2024treams} D.~Beutel, I.~Fernandez-Corbaton, and
C.~Rockstuhl, ``treams -- a T-matrix-based scattering code for
nanophotonics,'' Comput.\ Phys.\ Commun.\ \textbf{297}, 109076 (2024).
\bibitem{NAMM} W.~J.~Padilla, \emph{The Physics of the Neural Analytic
Modal Method} (internal document).
\end{thebibliography}

\end{document}
